%% ============================================================
%% sections/s360.tex
%% United States Code, Title 21 - § 360. Registration of producers of drugs or devices
%% (FD&C Act § 510). Source identifier: /us/usc/t21/s360
%% Relevance to the Physical AI oncology trial context:
%% Registration of producers; subsection (k) is the 510(k) premarket-notification pathway (distinct from § 360k state preemption).
%% Reproduced faithfully from the OLRC USLM XML (through Pub. L. 119-93).
%% No bracketed instructions or file names are inserted into the law.
%% ============================================================
\uscsection{§~360.}{Registration of producers of drugs or devices}
\uscprov{1}{\textbf{(a)}\textbf{ Definitions} As used in this section-}
\uscprov{2}{\textbf{(1)} the term ``manufacture, preparation, propagation, compounding, or processing'' shall include repackaging or otherwise changing the container, wrapper, or labeling of any drug package or device package in furtherance of the distribution of the drug or device from the original place of manufacture to the person who makes final delivery or sale to the ultimate consumer or user; and}
\uscprov{2}{\textbf{(2)} the term ``name'' shall include in the case of a partnership the name of each partner and, in the case of a corporation, the name of each corporate officer and director, and the State of incorporation.}
\uscprov{1}{\textbf{(b)}\textbf{ Annual registration}}
\uscprov{2}{\textbf{(1)} During the period beginning on October 1 and ending on December 31 of each year, every person who owns or operates any establishment in any State engaged in the manufacture, preparation, propagation, compounding, or processing of a drug or drugs shall register with the Secretary the name of such person, places of business of such person, all such establishments, the unique facility identifier of each such establishment, and a point of contact e-mail address.}
\uscprov{2}{\textbf{(2)} During the period beginning on October 1 and ending on December 31 of each year, every person who owns or operates any establishment in any State engaged in the manufacture, preparation, propagation, compounding, or processing of a device or devices shall register with the Secretary his name, places of business, and all such establishments.}
\uscprov{2}{\textbf{(3)} The Secretary shall specify the unique facility identifier system that shall be used by registrants under paragraph (1). The requirement to include a unique facility identifier in a registration under paragraph (1) shall not apply until the date that the identifier system is specified by the Secretary under the preceding sentence.}
\uscprov{1}{\textbf{(c)}\textbf{ New producers} Every person upon first engaging in the manufacture, preparation, propagation, compounding, or processing of a drug or drugs or a device or devices in any establishment which he owns or operates in any State shall immediately register with the Secretary-}
\uscprov{2}{\textbf{(1)} with respect to drugs, the information described under subsection (b)(1); and}
\uscprov{2}{\textbf{(2)} with respect to devices, the information described under subsection (b)(2)..\textsuperscript{1}\,{\footnotesize ~So in original.}}
\uscprov{1}{\textbf{(d)}\textbf{ Additional establishments} Every person duly registered in accordance with the foregoing subsections of this section shall immediately register with the Secretary any additional establishment which he owns or operates in any State and in which he begins the manufacture, preparation, propagation, compounding, or processing of a drug or drugs or a device or devices.}
\uscprov{1}{\textbf{(e)}\textbf{ Registration number; uniform system for identification of devices intended for human use} The Secretary may assign a registration number to any person or any establishment registered in accordance with this section. The Secretary may also assign a listing number to each drug or class of drugs listed under subsection (j). Any number assigned pursuant to the preceding sentence shall be the same as that assigned pursuant to the National Drug Code. The Secretary may by regulation prescribe a uniform system for the identification of devices intended for human use and may require that persons who are required to list such devices pursuant to subsection (j) shall list such devices in accordance with such system.}
\uscprov{1}{\textbf{(f)}\textbf{ Availability of registrations for inspection} The Secretary shall make available for inspection, to any person so requesting, any registration filed pursuant to this section; except that any list submitted pursuant to paragraph (3) of subsection (j) and the information accompanying any list or notice filed under paragraph (1) or (2) of that subsection shall be exempt from such inspection unless the Secretary finds that such an exemption would be inconsistent with protection of the public health.}
\uscprov{1}{\textbf{(g)}\textbf{ Exclusions from application of section} The foregoing subsections of this section shall not apply to-}
\uscprov{2}{\textbf{(1)} pharmacies which maintain establishments in conformance with any applicable local laws regulating the practice of pharmacy and medicine and which are regularly engaged in dispensing prescription drugs or devices, upon prescriptions of practitioners licensed to administer such drugs or devices to patients under the care of such practitioners in the course of their professional practice, and which do not manufacture, prepare, propagate, compound, or process drugs or devices for sale other than in the regular course of their business of dispensing or selling drugs or devices at retail;}
\uscprov{2}{\textbf{(2)} practitioners licensed by law to prescribe or administer drugs or devices and who manufacture, prepare, propagate, compound, or process drugs or devices solely for use in the course of their professional practice;}
\uscprov{2}{\textbf{(3)} persons who manufacture, prepare, propagate, compound, or process drugs or devices solely for use in research, teaching, or chemical analysis and not for sale;}
\uscprov{2}{\textbf{(4)} any distributor who acts as a wholesale distributor of devices, and who does not manufacture, repackage, process, or relabel a device; or}
\uscprov{2}{\textbf{(5)} such other classes of persons as the Secretary may by regulation exempt from the application of this section upon a finding that registration by such classes of persons in accordance with this section is not necessary for the protection of the public health.}
\uscprov{1}{In this subsection, the term ``wholesale distributor'' means any person (other than the manufacturer or the initial importer) who distributes a device from the original place of manufacture to the person who makes the final delivery or sale of the device to the ultimate consumer or user.}
\uscprov{1}{\textbf{(h)}\textbf{ Inspections}}
\uscprov{2}{\textbf{(1)}\textbf{ In general} Every establishment that is required to be registered with the Secretary under this section shall be subject to inspection pursuant to section 374 of this title.}
\uscprov{2}{\textbf{(2)}\textbf{ Risk-based schedule for devices}}
\uscprov{3}{\textbf{(A)}\textbf{ In general} The Secretary, acting through one or more officers or employees duly designated by the Secretary, shall inspect establishments described in paragraph (1) that are engaged in the manufacture, propagation, compounding, or processing of a device or devices (referred to in this subsection as ``device establishments'') in accordance with a risk-based schedule established by the Secretary.}
\uscprov{3}{\textbf{(B)}\textbf{ Factors and considerations} In establishing the risk-based schedule under subparagraph (A), the Secretary shall-}
\uscprov{4}{\textbf{(i)} apply, to the extent applicable for device establishments, the factors identified in paragraph (4); and}
\uscprov{4}{\textbf{(ii)} consider the participation of the device establishment, as applicable, in international device audit programs in which the United States participates or the United States recognizes for purposes of inspecting device establishments.}
\uscprov{2}{\textbf{(3)}\textbf{ Risk-based schedule for drugs} The Secretary, acting through one or more officers or employees duly designated by the Secretary, shall inspect establishments described in paragraph (1) that are engaged in the manufacture, preparation, propagation, compounding, or processing of a drug or drugs (referred to in this subsection as ``drug establishments'') in accordance with a risk-based schedule established by the Secretary.}
\uscprov{2}{\textbf{(4)}\textbf{ Risk factors} In establishing a risk-based schedule under paragraph (2) or (3), the Secretary shall inspect establishments according to the known safety risks of such establishments, which shall be based on the following factors:}
\uscprov{3}{\textbf{(A)} The compliance history of the establishment.}
\uscprov{3}{\textbf{(B)} The record, history, and nature of recalls linked to the establishment.}
\uscprov{3}{\textbf{(C)} The inherent risk of the drug or device manufactured, prepared, propagated, compounded, or processed at the establishment.}
\uscprov{3}{\textbf{(D)} The inspection frequency and history of the establishment, including whether the establishment has been inspected pursuant to section 374 of this title within the last 4 years.}
\uscprov{3}{\textbf{(E)} Whether the establishment has been inspected by a foreign government or an agency of a foreign government recognized under section 384e of this title.}
\uscprov{3}{\textbf{(F)} The compliance history of establishments in the country or region in which the establishment is located that are subject to regulation under this chapter, including the history of violations related to products exported from such country or region that are subject to such regulation.}
\uscprov{3}{\textbf{(G)} Any other criteria deemed necessary and appropriate by the Secretary for purposes of allocating inspection resources.}
\uscprov{2}{\textbf{(5)}\textbf{ Effect of status} In determining the risk associated with an establishment for purposes of establishing a risk-based schedule under paragraph (3), the Secretary shall not consider whether the drugs manufactured, prepared, propagated, compounded, or processed by such establishment are drugs described in section 353(b) of this title.}
\uscprov{2}{\textbf{(6)}\textbf{ Annual report on inspections of establishments} Not later than May 1 of each year, the Secretary shall make available on the Internet Web site of the Food and Drug Administration a report regarding-}
\uscprov{3}{\textbf{(A)}}
\uscprov{4}{\textbf{(i)} the number of domestic and foreign establishments registered pursuant to this section in the previous fiscal year;}
\uscprov{4}{\textbf{(ii)} the number of such registered establishments in each region of interest;}
\uscprov{4}{\textbf{(iii)} the number of such domestic establishments and the number of such foreign establishments, including the number of establishments in each region of interest, that the Secretary inspected in the previous fiscal year;}
\uscprov{4}{\textbf{(iv)} the number of inspections to support actions by the Secretary on applications under section 355 of this title or section 262 of title 42, including the number of inspections to support actions by the Secretary on supplemental applications, including changes to manufacturing processes, the Secretary conducted in the previous fiscal year;}
\uscprov{4}{\textbf{(v)} the number of routine surveillance inspections the Secretary conducted in the previous fiscal year, including in each region of interest;}
\uscprov{4}{\textbf{(vi)} the number of for-cause inspections the Secretary conducted in the previous fiscal year, not including inspections described in clause (iv), including in each region of interest; and}
\uscprov{4}{\textbf{(vii)} the number of inspections the Secretary has recognized pursuant to an agreement entered into pursuant to section 384e of this title, or otherwise recognized, for each of the types of inspections described in clauses (v) and (vi), including for inspections of establishments in each region of interest.\textsuperscript{2}\,{\footnotesize ~So in original. The period probably should be a semicolon.}}
\uscprov{3}{\textbf{(B)} with respect to establishments that manufacture, prepare, propagate, compound, or process an active ingredient of a drug or a finished drug product, the number of each such type of establishment;}
\uscprov{3}{\textbf{(C)} the percentage of the budget of the Food and Drug Administration used to fund the inspections described under subparagraph (A); and}
\uscprov{3}{\textbf{(D)} the status of the efforts of the Food and Drug Administration to expand its recognition of inspections conducted or recognized by foreign regulatory authorities under section 384e of this title, including any obstacles to expanding the use of such recognition.}
\uscprov{2}{\textbf{(7)}\textbf{ Region of interest} For purposes of paragraph (6)(A), the term ``region of interest'' means a foreign geographic region or country, including the People's Republic of China, India, the European Union, the United Kingdom, and any other country or geographic region, as the Secretary determines appropriate.}
\uscprov{1}{\textbf{(i)}\textbf{ Registration of foreign establishments}}
\uscprov{2}{\textbf{(1)} Every person who owns or operates any establishment within any foreign country engaged in the manufacture, preparation, propagation, compounding, or processing of a drug or device that is imported or offered for import into the United States shall, through electronic means in accordance with the criteria of the Secretary-}
\uscprov{3}{\textbf{(A)} upon first engaging in any such activity, immediately submit a registration to the Secretary that includes-}
\uscprov{4}{\textbf{(i)} with respect to drugs, the name and place of business of such person, all such establishments, the unique facility identifier of each such establishment, a point of contact e-mail address, the name of the United States agent of each such establishment, the name of each importer of such drug in the United States that is known to the establishment, and the name of each person who imports or offers for import such drug to the United States for purposes of importation; and}
\uscprov{4}{\textbf{(ii)} with respect to devices, the name and place of business of the establishment, the name of the United States agent for the establishment, the name of each importer of such device in the United States that is known to the establishment, and the name of each person who imports or offers for import such device to the United States for purposes of importation; and}
\uscprov{3}{\textbf{(B)} each establishment subject to the requirements of subparagraph (A) shall thereafter register with the Secretary during the period beginning on October 1 and ending on December 31 of each year.}
\uscprov{2}{\textbf{(2)} The establishment shall also provide the information required by subsection (j).}
\uscprov{2}{\textbf{(3)} The Secretary is authorized to enter into cooperative arrangements with officials of foreign countries to ensure that adequate and effective means are available for purposes of determining, from time to time, whether drugs or devices manufactured, prepared, propagated, compounded, or processed by an establishment described in paragraph (1), if imported or offered for import into the United States, shall be refused admission on any of the grounds set forth in section 381(a) of this title.}
\uscprov{2}{\textbf{(4)} The Secretary shall specify the unique facility identifier system that shall be used by registrants under paragraph (1) with respect to drugs. The requirement to include a unique facility identifier in a registration under paragraph (1) with respect to drugs shall not apply until the date that the identifier system is specified by the Secretary under the preceding sentence.}
\uscprov{2}{\textbf{(5)} The requirements of paragraphs (1) and (2) shall apply regardless of whether the drug or device undergoes further manufacture, preparation, propagation, compounding, or processing at a separate establishment outside the United States prior to being imported or offered for import into the United States.}
\uscprov{1}{\textbf{(j)}\textbf{ Filing of lists of drugs and devices manufactured, prepared, propagated and compounded by registrants; statements; accompanying disclosures}}
\uscprov{2}{\textbf{(1)} Every person who registers with the Secretary under subsection (b), (c), (d), or (i) shall, at the time of registration under any such subsection, file with the Secretary a list of all drugs and a list of all devices and a brief statement of the basis for believing that each device included in the list is a device rather than a drug (with each drug and device in each list listed by its established name (as defined in section 352(e) of this title) and by any proprietary name) which are being manufactured, prepared, propagated, compounded, or processed by him for commercial distribution and which he has not included in any list of drugs or devices filed by him with the Secretary under this paragraph or paragraph (2) before such time of registration. Such list shall be prepared in such form and manner as the Secretary may prescribe and shall be accompanied by-}
\uscprov{3}{\textbf{(A)} in the case of a drug contained in the applicable list and subject to section 355 or 360b of this title, or a device intended for human use contained in the applicable list with respect to which a performance standard has been established under section 360d of this title or which is subject to section 360e of this title, a reference to the authority for the marketing of such drug or device and a copy of all labeling for such drug or device;}
\uscprov{3}{\textbf{(B)} in the case of any other drug or device contained in an applicable list-}
\uscprov{4}{\textbf{(i)} which drug is subject to section 353(b)(1) of this title, or which device is a restricted device, a copy of all labeling for such drug or device, a representative sampling of advertisements for such drug or device, and, upon request made by the Secretary for good cause, a copy of all advertisements for a particular drug product or device, or}
\uscprov{4}{\textbf{(ii)} which drug is not subject to section 353(b)(1) of this title or which device is not a restricted device, the label and package insert for such drug or device and a representative sampling of any other labeling for such drug or device;}
\uscprov{3}{\textbf{(C)} in the case of any drug contained in an applicable list which is described in subparagraph (B), a quantitative listing of its active ingredient or ingredients, except that with respect to a particular drug product the Secretary may require the submission of a quantitative listing of all ingredients if he finds that such submission is necessary to carry out the purposes of this chapter;}
\uscprov{3}{\textbf{(D)} if the registrant filing a list has determined that a particular drug product or device contained in such list is not subject to section 355 or 360b of this title, or the particular device contained in such list is not subject to a performance standard established under section 360d of this title or to section 360e of this title or is not a restricted device a brief statement of the basis upon which the registrant made such determination if the Secretary requests such a statement with respect to that particular drug product or device; and}
\uscprov{3}{\textbf{(E)} in the case of a drug contained in the applicable list, the name and place of business of each manufacturer of an excipient of the listed drug with which the person listing the drug conducts business, including all establishments used in the production of such excipient, the unique facility identifier of each such establishment, and a point of contact e-mail address for each such excipient manufacturer.}
\uscprov{2}{\textbf{(2)} Each person who registers with the Secretary under this section shall report to the Secretary, with regard to drugs once during the month of June of each year and once during the month of December of each year, and with regard to devices once each year during the period beginning on October 1 and ending on December 31, the following information:}
\uscprov{3}{\textbf{(A)} A list of each drug or device introduced by the registrant for commercial distribution which has not been included in any list previously filed by him with the Secretary under this subparagraph or paragraph (1) of this subsection. A list under this subparagraph shall list a drug or device by its established name (as defined in section 352(e) of this title), and by any proprietary name it may have and shall be accompanied by the other information required by paragraph (1).}
\uscprov{3}{\textbf{(B)} If since the date the registrant last made a report under this paragraph (or if he has not made a report under this paragraph, since February 1, 1973) he has discontinued the manufacture, preparation, propagation, compounding, or processing for commercial distribution of a drug or device included in a list filed by him under subparagraph (A) or paragraph (1); notice of such discontinuance, the date of such discontinuance, and the identity (by established name (as defined in section 352(e) of this title) and by any proprietary name) of such drug or device.}
\uscprov{3}{\textbf{(C)} If since the date the registrant reported pursuant to subparagraph (B) a notice of discontinuance he has resumed the manufacture, preparation, propagation, compounding, or processing for commercial distribution of the drug or device with respect to which such notice of discontinuance was reported; notice of such resumption, the date of such resumption, the identity of such drug or device (each by established name (as defined in section 352(e) of this title) and by any proprietary name), and the other information required by paragraph (1), unless the registrant has previously reported such resumption to the Secretary pursuant to this subparagraph.}
\uscprov{3}{\textbf{(D)} Any material change in any information previously submitted pursuant to this paragraph or paragraph (1).}
\uscprov{2}{\textbf{(3)}}
\uscprov{3}{\textbf{(A)} Each person who registers with the Secretary under this section with regard to a drug shall report annually to the Secretary on the amount of each drug listed under paragraph (1) that was manufactured, prepared, propagated, compounded, or processed by such person for commercial distribution. Such information may be required to be submitted in an electronic format as determined by the Secretary. The Secretary may require that information required to be reported under this paragraph be submitted at the time a public health emergency is declared by the Secretary under section 247d of title 42.}
\uscprov{3}{\textbf{(B)} By order of the Secretary, certain biological products or categories of biological products regulated under section 262 of title 42 may be exempt from some or all of the reporting requirements under subparagraph (A), if the Secretary determines that applying such reporting requirements to such biological products or categories of biological products is not necessary to protect the public health.}
\uscprov{2}{\textbf{(4)} The Secretary may also require each registrant under this section to submit a list of each drug product which (A) the registrant is manufacturing, preparing, propagating, compounding, or processing for commercial distribution, and (B) contains a particular ingredient. The Secretary may not require the submission of such a list unless he has made a finding that the submission of such a list is necessary to carry out the purposes of this chapter.}
\uscprov{2}{\textbf{(5)} The Secretary shall require persons subject to this subsection to use, for purposes of this subsection, the unique facility identifier systems specified under subsections (b)(3) and (i)(4) with respect to drugs. Such requirement shall not apply until the date that the identifier system under subsection (b)(3) or (i)(4), as applicable, is specified by the Secretary.}
\uscprov{1}{\textbf{(k)}\textbf{ Report preceding introduction of devices into interstate commerce} Each person who is required to register under this section and who proposes to begin the introduction or delivery for introduction into interstate commerce for commercial distribution of a device intended for human use shall, at least ninety days before making such introduction or delivery, report to the Secretary or person who is accredited under section 360m(a) of this title (in such form and manner as the Secretary shall by regulation prescribe)-}
\uscprov{2}{\textbf{(1)} the class in which the device is classified under section 360c of this title or if such person determines that the device is not classified under such section, a statement of that determination and the basis for such person's determination that the device is or is not so classified, and}
\uscprov{2}{\textbf{(2)} action taken by such person to comply with requirements under section 360d or 360e of this title which are applicable to the device.}
\uscprov{1}{A notification submitted under this subsection that contains clinical trial data for an applicable device clinical trial (as defined in section 282(j)(1) of title 42) shall be accompanied by the certification required under section 282(j)(5)(B) of such title. Such certification shall not be considered an element of such notification.}
\uscprov{1}{\textbf{(l)}\textbf{ Exemption from reporting requirements}}
\uscprov{2}{\textbf{(1)} A report under subsection (k) is not required for a device intended for human use that is exempted from the requirements of this subsection under subsection (m) or is within a type that has been classified into class I under section 360c of this title, or with respect to a change that is consistent with a predetermined change control plan cleared under section 360e-4 of this title. The exception established in the preceding sentence does not apply to any class I device that is intended for a use which is of substantial importance in preventing impairment of human health, or to any class I device that presents a potential unreasonable risk of illness or injury.}
\uscprov{2}{\textbf{(2)} Not later than 120 calendar days after December 13, 2016, and at least once every 5 years thereafter, as the Secretary determines appropriate, the Secretary shall identify, through publication in the Federal Register, any type of class I device that the Secretary determines no longer requires a report under subsection (k) to provide reasonable assurance of safety and effectiveness. Upon such publication-}
\uscprov{3}{\textbf{(A)} each type of class I device so identified shall be exempt from the requirement for a report under subsection (k); and}
\uscprov{3}{\textbf{(B)} the classification regulation applicable to each such type of device shall be deemed amended to incorporate such exemption.}
\uscprov{1}{\textbf{(m)}\textbf{ List of exempt class II devices; initial and final determinations by Secretary; publication in Federal Register}}
\uscprov{2}{\textbf{(1)} The Secretary shall-}
\uscprov{3}{\textbf{(A)} not later than 90 days after December 13, 2016, and at least once every 5 years thereafter, as the Secretary determines appropriate-}
\uscprov{4}{\textbf{(i)} publish in the Federal Register a notice that contains a list of each type of class II device that the Secretary determines no longer requires a report under subsection (k) to provide reasonable assurance of safety and effectiveness; and}
\uscprov{4}{\textbf{(ii)} provide for a period of not less than 60 calendar days for public comment beginning on the date of the publication of such notice; and}
\uscprov{3}{\textbf{(B)} not later than 210 calendar days after December 13, 2016, publish in the Federal Register a list representing the Secretary's final determination with respect to the devices contained in the list published under subparagraph (A).}
\uscprov{2}{\textbf{(2)} Beginning on the date that is 1 calendar day after the date of publication of the final list under paragraph (1)(B), the Secretary may exempt a class II device from the requirement to submit a report under subsection (k), upon the Secretary's own initiative or a petition of an interested person, if the Secretary determines that such report is not necessary to assure the safety and effectiveness of the device. The Secretary shall publish in the Federal Register notice of the intent of the Secretary to exempt the device, or of the petition, and provide a 60-calendar-day period for public comment. Within 120 days after the issuance of the notice in the Federal Register, the Secretary shall publish an order in the Federal Register that sets forth the final determination of the Secretary regarding the exemption of the device that was the subject of the notice. If the Secretary fails to respond to a petition within 180 days of receiving it, the petition shall be deemed to be granted.}
\uscprov{2}{\textbf{(3)} Upon the publication of the final list under paragraph (1)(B)-}
\uscprov{3}{\textbf{(A)} each type of class II device so listed shall be exempt from the requirement for a report under subsection (k); and}
\uscprov{3}{\textbf{(B)} the classification regulation applicable to each such type of device shall be deemed amended to incorporate such exemption.}
\uscprov{1}{\textbf{(n)}\textbf{ Review of report; time for determination by Secretary}}
\uscprov{2}{\textbf{(1)} The Secretary shall review the report required in subsection (k) and make a determination under section 360c(f)(1) of this title not later than 90 days after receiving the report.}
\uscprov{2}{\textbf{(2)}}
\uscprov{3}{\textbf{(A)} Not later than 18 months after July 9, 2012, the Secretary shall submit to the Committee on Energy and Commerce of the House of Representatives and the Committee on Health, Education, Labor, and Pensions of the Senate a report regarding when a premarket notification under subsection (k) should be submitted for a modification or change to a legally marketed device. The report shall include the Secretary's interpretation of the following terms: ``could significantly affect the safety or effectiveness of the device'', ``a significant change or modification in design, material, chemical composition, energy source, or manufacturing process'', and ``major change or modification in the intended use of the device''. The report also shall discuss possible processes for industry to use to determine whether a new submission under subsection (k) is required and shall analyze how to leverage existing quality system requirements to reduce premarket burden, facilitate continual device improvement, and provide reasonable assurance of safety and effectiveness of modified devices. In developing such report, the Secretary shall consider the input of interested stakeholders.}
\uscprov{3}{\textbf{(B)} The Secretary shall withdraw the Food and Drug Administration draft guidance entitled ``Guidance for Industry and FDA Staff-510(k) Device Modifications: Deciding When to Submit a 510(k) for a Change to an Existing Device'', dated July 27, 2011, and shall not use this draft guidance as part of, or for the basis of, any premarket review or any compliance or enforcement decisions or actions. The Secretary shall not issue-}
\uscprov{4}{\textbf{(i)} any draft guidance or proposed regulation that addresses when to submit a premarket notification submission for changes and modifications made to a manufacturer's previously cleared device before the receipt by the Committee on Energy and Commerce of the House of Representatives and the Committee on Health, Education, Labor, and Pensions of the Senate of the report required in subparagraph (A); and}
\uscprov{4}{\textbf{(ii)} any final guidance or regulation on that topic for one year after date of receipt of such report by the Committee on Energy and Commerce of the House of Representatives and the Committee on Health, Education, Labor, and Pensions of the Senate.}
\uscprov{3}{\textbf{(C)} The Food and Drug Administration guidance entitled ``Deciding When to Submit a 510(k) for a Change to an Existing Device'', dated January 10, 1997, shall be in effect until the subsequent issuance of guidance or promulgation, if appropriate, of a regulation described in subparagraph (B), and the Secretary shall interpret such guidance in a manner that is consistent with the manner in which the Secretary has interpreted such guidance since 1997.}
\uscprov{1}{\textbf{(o)}\textbf{ Reprocessed single-use devices}}
\uscprov{2}{\textbf{(1)} With respect to reprocessed single-use devices for which reports are required under subsection (k):}
\uscprov{3}{\textbf{(A)} The Secretary shall identify such devices or types of devices for which reports under such subsection must, in order to ensure that the device is substantially equivalent to a predicate device, include validation data, the types of which shall be specified by the Secretary, regarding cleaning and sterilization, and functional performance demonstrating that the single-use device will remain substantially equivalent to its predicate device after the maximum number of times the device is reprocessed as intended by the person submitting the premarket notification. Within six months after October 26, 2002, the Secretary shall publish in the Federal Register a list of the types so identified, and shall revise the list as appropriate. Reports under subsection (k) for devices or types of devices within a type included on the list are, upon publication of the list, required to include such validation data.}
\uscprov{3}{\textbf{(B)} In the case of each report under subsection (k) that was submitted to the Secretary before the publication of the initial list under subparagraph (A), or any revision thereof, and was for a device or type of device included on such list, the person who submitted the report under subsection (k) shall submit validation data as described in subparagraph (A) to the Secretary not later than nine months after the publication of the list. During such nine-month period, the Secretary may not take any action under this chapter against such device solely on the basis that the validation data for the device have not been submitted to the Secretary. After the submission of the validation data to the Secretary, the Secretary may not determine that the device is misbranded under section 352(\textit{o}) of this title or adulterated under section 351(f)(1)(B) of this title, or take action against the device under section 331(p) of this title for failure to provide any information required by subsection (k) until (i) the review is terminated by withdrawal of the submission of the report under subsection (k); (ii) the Secretary finds the data to be acceptable and issues a letter; or (iii) the Secretary determines that the device is not substantially equivalent to a predicate device. Upon a determination that a device is not substantially equivalent to a predicate device, or if such submission is withdrawn, the device can no longer be legally marketed.}
\uscprov{3}{\textbf{(C)} In the case of a report under subsection (k) for a device identified under subparagraph (A) that is of a type for which the Secretary has not previously received a report under such subsection, the Secretary may, in advance of revising the list under subparagraph (A) to include such type, require that the report include the validation data specified in subparagraph (A).}
\uscprov{3}{\textbf{(D)} Section 352(\textit{o}) of this title applies with respect to the failure of a report under subsection (k) to include validation data required under subparagraph (A).}
\uscprov{2}{\textbf{(2)} With respect to critical or semi-critical reprocessed single-use devices that, under subsection (\textit{l}) or (m), are exempt from the requirement of submitting reports under subsection (k):}
\uscprov{3}{\textbf{(A)} The Secretary shall identify such devices or types of devices for which such exemptions should be terminated in order to provide a reasonable assurance of the safety and effectiveness of the devices. The Secretary shall publish in the Federal Register a list of the devices or types of devices so identified, and shall revise the list as appropriate. The exemption for each device or type included on the list is terminated upon the publication of the list. For each report under subsection (k) submitted pursuant to this subparagraph the Secretary shall require the validation data described in paragraph (1)(A).}
\uscprov{3}{\textbf{(B)} For each device or type of device included on the list under subparagraph (A), a report under subsection (k) shall be submitted to the Secretary not later than 15 months after the publication of the initial list, or a revision of the list, whichever terminates the exemption for the device. During such 15-month period, the Secretary may not take any action under this chapter against such device solely on the basis that such report has not been submitted to the Secretary. After the submission of the report to the Secretary the Secretary may not determine that the device is misbranded under section 352(\textit{o}) of this title or adulterated under section 351(f)(1)(B) of this title, or take action against the device under section 331(p) of this title for failure to provide any information required by subsection (k) until (i) the review is terminated by withdrawal of the submission; (ii) the Secretary determines by order that the device is substantially equivalent to a predicate device; or (iii) the Secretary determines by order that the device is not substantially equivalent to a predicate device. Upon a determination that a device is not substantially equivalent to a predicate device, the device can no longer be legally marketed.}
\uscprov{3}{\textbf{(C)} In the case of semi-critical devices, the initial list under subparagraph (A) shall be published not later than 18 months after the effective date of this subsection. In the case of critical devices, the initial list under such subparagraph shall be published not later than six months after such effective date.}
\uscprov{3}{\textbf{(D)} Section 352(\textit{o}) of this title applies with respect to the failure to submit a report under subsection (k) that is required pursuant to subparagraph (A), including a failure of the report to include validation data required in such subparagraph.}
\uscprov{3}{\textbf{(E)} The termination under subparagraph (A) of an exemption under subsection (\textit{l}) or (m) for a critical or semi-critical reprocessed single-use device does not terminate the exemption under subsection (\textit{l}) or (m) for the original device.}
\uscprov{1}{\textbf{(p)}\textbf{ Electronic registration and listing}}
\uscprov{2}{\textbf{(1)}\textbf{ In general} Registrations and listings under this section (including the submission of updated information) shall be submitted to the Secretary by electronic means unless the Secretary grants a request for waiver of such requirement because use of electronic means is not reasonable for the person requesting such waiver.}
\uscprov{2}{\textbf{(2)}\textbf{ Electronic database} Not later than 2 years after the Secretary specifies a unique facility identifier system under subsections (b) and (i), the Secretary shall maintain an electronic database, which shall not be subject to inspection under subsection (f), populated with the information submitted as described under paragraph (1) that-}
\uscprov{3}{\textbf{(A)} enables personnel of the Food and Drug Administration to search the database by any field of information submitted in a registration described under paragraph (1), or combination of such fields; and}
\uscprov{3}{\textbf{(B)} uses the unique facility identifier system to link with other relevant databases within the Food and Drug Administration, including the database for submission of information under section 381(r) of this title.}
\uscprov{2}{\textbf{(3)}\textbf{ Risk-based information and coordination} The Secretary shall ensure the accuracy and coordination of relevant Food and Drug Administration databases in order to identify and inform risk-based inspections under subsection (h).}
\uscprov{1}{\textbf{(q)}\textbf{ Reusable medical devices}}
\uscprov{2}{\textbf{(1)}\textbf{ In general} Not later than 180 days after December 13, 2016, the Secretary shall identify and publish a list of reusable device types for which reports under subsection (k) are required to include-}
\uscprov{3}{\textbf{(A)} instructions for use, which have been validated in a manner specified by the Secretary; and}
\uscprov{3}{\textbf{(B)} validation data, the types of which shall be specified by the Secretary;}
\uscprov{2}{regarding cleaning, disinfection, and sterilization, and for which a substantial equivalence determination may be based.}
\uscprov{2}{\textbf{(2)}\textbf{ Revision of list} The Secretary shall revise the list under paragraph (2),\textsuperscript{3}\,{\footnotesize ~So in original. Probably should be ``paragraph (1),''.} as the Secretary determines appropriate, with notice in the Federal Register.}
\uscprov{2}{\textbf{(3)}\textbf{ Content of reports} Reports under subsection (k) that are submitted after the publication of the list described in paragraph (1), for devices or types of devices included on such list, shall include such instructions for use and validation data.}
\uscsource{(June 25, 1938, ch. 675, §~510, as added Pub. L. 87-781, title III, §~302, Oct. 10, 1962, 76 Stat. 794; amended Pub. L. 89-74, §~4, July 15, 1965, 79 Stat. 231; Pub. L. 91-513, title II, §~701(e), Oct. 27, 1970, 84 Stat. 1282; Pub. L. 92-387, §§~3, 4(a)-(c), Aug. 16, 1972, 86 Stat. 560-562; Pub. L. 94-295, §~4(a), May 28, 1976, 90 Stat. 579; Pub. L. 105-115, title I, §~125(a)(2)(C), title II, §§~206(a), 209(a), 213(b), title IV, §~417, Nov. 21, 1997, 111 Stat. 2325, 2338, 2341, 2347, 2379; Pub. L. 107-188, title III, §~321(a), June 12, 2002, 116 Stat. 675; Pub. L. 107-250, title II, §§~201(e), 207, 211, title III, §~302(b), Oct. 26, 2002, 116 Stat. 1609, 1613, 1614, 1616; Pub. L. 108-214, §~2(c)(2), Apr. 1, 2004, 118 Stat. 576; Pub. L. 110-85, title II, §§~222-224, title VIII, §~801(b)(3)(C), Sept. 27, 2007, 121 Stat. 853, 921; Pub. L. 112-144, title VI, §~604, title VII, §§~701, 702(b)-705, July 9, 2012, 126 Stat. 1052, 1064-1066; Pub. L. 114-255, div. A, title III, §§~3054, 3059(a), 3101(a)(2)(H), Dec. 13, 2016, 130 Stat. 1126, 1130, 1154; Pub. L. 115-52, title VII, §~701(a), title IX, §~901(e), Aug. 18, 2017, 131 Stat. 1054, 1076; Pub. L. 116-136, div. A, title III, §~3112(e), Mar. 27, 2020, 134 Stat. 363; Pub. L. 117-328, div. FF, title II, §§~2511(a), 2515(a)(3), title III, §§~3308(b)(1), 3613(a), 3616(c), Dec. 29, 2022, 136 Stat. 5803, 5806, 5836, 5872, 5875.)}
\usccrosshead{Editorial Notes}
\uscnotehead{References in Text}
\uscnotepar{The effective date of this subsection, referred to in subsec. (\textit{o})(2)(C), probably means the date of the enactment of Pub. L. 107-250, which enacted subsec. (\textit{o}) of this section and was approved Oct. 26, 2002.}
\uscnotehead{Amendments}
\uscnotepar{2022-Subsec. (h)(4)(F), (G). Pub. L. 117-328, §~3613(a), added subpar. (F) and redesignated former subpar. (F) as (G).}
\uscnotepar{Subsec. (h)(6). Pub. L. 117-328, §~3616(c)(1)(A), substituted ``Not'' for ``Beginning in 2014, not'' in introductory provisions.}
\uscnotepar{Subsec. (h)(6)(A). Pub. L. 117-328, §~3616(c)(1)(B), amended subpar. (A) generally. Prior to amendment, subpar. (A) read as follows:}
\uscnotepar{``(A)(i) the number of domestic and foreign establishments registered pursuant to this section in the previous calendar year; and}
\uscnotepar{``(ii) the number of such domestic establishments and the number of such foreign establishments that the Secretary inspected in the previous calendar year;''.}
\uscnotepar{Subsec. (h)(6)(D). Pub. L. 117-328, §~3616(c)(1)(C)-(E), added subpar. (D).}
\uscnotepar{Subsec. (h)(7). Pub. L. 117-328, §~3616(c)(2), added par. (7).}
\uscnotepar{Subsec. (i)(5). Pub. L. 117-328, §~2511(a), added par. (5).}
\uscnotepar{Subsec. (j)(3) to (5). Pub. L. 117-328, §~2515(a)(3), made technical amendment to directory language of Pub. L. 116-136, §~3112(e). See 2020 Amendment note below.}
\uscnotepar{Subsec. (\textit{l})(1). Pub. L. 117-328, §~3308(b)(1), inserted ``, or with respect to a change that is consistent with a predetermined change control plan cleared under section 360e-4 of this title'' after ``section 360c of this title''.}
\uscnotepar{2020-Subsec. (j)(3) to (5). Pub. L. 116-136, §~3112(e), as amended by Pub. L. 117-328, §~2515(a)(3), added par. (3) and redesignated former pars. (3) and (4) as (4) and (5), respectively.}
\uscnotepar{2017-Subsec. (h)(2). Pub. L. 115-52, §~701(a)(1), added par. (2) and struck out former par. (2). Prior to amendment, text read as follows: ``Every establishment described in paragraph (1), in any State, that is engaged in the manufacture, propagation, compounding, or processing of a device or devices classified in class II or III shall be so inspected by one or more officers or employees duly designated by the Secretary, or by persons accredited to conduct inspections under section 374(g) of this title, at least once in the 2-year period beginning with the date of registration of such establishment pursuant to this section and at least once in every successive 2-year period thereafter.''}
\uscnotepar{Subsec. (h)(4). Pub. L. 115-52, §~701(a)(2)(A), substituted ``paragraph (2) or (3)'' for ``paragraph (3)'' in introductory provisions.}
\uscnotepar{Subsec. (h)(4)(C). Pub. L. 115-52, §~701(a)(2)(B), inserted ``or device'' after ``drug''.}
\uscnotepar{Subsec. (h)(6). Pub. L. 115-52, §~901(e), substituted ``May 1'' for ``February 1'' in introductory provisions.}
\uscnotepar{2016-Subsec. (h)(4). Pub. L. 114-255, §~3101(a)(2)(H)(i), substituted ``establishing a risk-based schedule'' for ``establishing the risk-based scheduled'' in introductory provisions.}
\uscnotepar{Subsec. (h)(6)(A). Pub. L. 114-255, §~3101(a)(2)(H)(ii)(I), substituted ``calendar'' for ``fiscal'' in cls. (i) and (ii).}
\uscnotepar{Subsec. (h)(6)(B). Pub. L. 114-255, §~3101(a)(2)(H)(ii)(II), substituted ``an active ingredient of a drug or a finished drug product'' for ``an active ingredient of a drug, a finished drug product, or an excipient of a drug''.}
\uscnotepar{Subsec. (\textit{l}). Pub. L. 114-255, §~3054(a), designated existing provisions as par. (1) and added par. (2).}
\uscnotepar{Subsec. (m)(1). Pub. L. 114-255, §~3054(b)(1), added par. (1) and struck out former par. (1) which read as follows: ``Not later than 60 days after November 21, 1997, the Secretary shall publish in the Federal Register a list of each type of class II device that does not require a report under subsection (k) to provide reasonable assurance of safety and effectiveness. Each type of class II device identified by the Secretary as not requiring the report shall be exempt from the requirement to provide a report under subsection (k) as of the date of the publication of the list in the Federal Register. The Secretary shall publish such list on the Internet site of the Food and Drug Administration. The list so published shall be updated not later than 30 days after each revision of the list by the Secretary.''}
\uscnotepar{Subsec. (m)(2). Pub. L. 114-255, §~3054(b)(2)(B), substituted ``60-calendar-day period'' for ``30-day period''.}
\uscnotepar{Pub. L. 114-255, §~3054(b)(2)(A), which directed the substitution of ``1 calendar day after the date of publication of the final list under paragraph (1)(B),'' for ``1 day after the date of publication of a list under this subsection,'', was executed by making the substitution for ``1 day after the date of the publication of a list under this subsection,'' to reflect the probable intent of Congress.}
\uscnotepar{Subsec. (m)(3). Pub. L. 114-255, §~3054(b)(2)(C), added par. (3).}
\uscnotepar{Subsec. (q). Pub. L. 114-255, §~3059(a), added subsec. (q).}
\uscnotepar{2012-Subsec. (b)(1). Pub. L. 112-144, §~701(1)(A), which directed amendment of par. (1) by ``striking `On or before' and all that follows through the period at the end and inserting the following: `During the period beginning on October 1 and ending on December 31 of each year, every person who owns or operates any establishment in any State engaged in the manufacture, preparation, propagation, compounding, or processing of a drug or drugs shall register with the Secretary the name of such person, places of business of such person, all such establishments, the unique facility identifier of each such establishment, and a point of contact e-mail address.; and'', was executed as if an end quotation mark for the inserted material followed ``address.'', to reflect the probable intent of Congress. Prior to amendment, stricken text read as follows: ``On or before December 31 of each year every person who owns or operates any establishment in any State engaged in the manufacture, preparation, propagation, compounding, or processing of a drug or drugs shall register with the Secretary his name, places of business, and all such establishments.''}
\uscnotepar{Subsec. (b)(3). Pub. L. 112-144, §~701(1)(B), added par. (3).}
\uscnotepar{Subsec. (c). Pub. L. 112-144, §~701(2), substituted ``with the Secretary-'' and pars. (1) and (2) for ``with the Secretary his name, place of business, and such establishment''.}
\uscnotepar{Subsec. (h). Pub. L. 112-144, §~705, amended subsec. (h) generally. Prior to amendment, text read as follows: ``Every establishment in any State registered with the Secretary pursuant to this section shall be subject to inspection pursuant to section 374 of this title and every such establishment engaged in the manufacture, propagation, compounding, or processing of a drug or drugs or of a device or devices classified in class II or III shall be so inspected by one or more officers or employees duly designated by the Secretary, or by persons accredited to conduct inspections under section 374(g) of this title, at least once in the two-year period beginning with the date of registration of such establishment pursuant to this section and at least once in every successive two-year period thereafter.''}
\uscnotepar{Subsec. (i)(1). Pub. L. 112-144, §~702(b)(1)(A), amended introductory provisions generally. Prior to amendment, text read as follows: ``Any establishment within any foreign country engaged in the manufacture, preparation, propagation, compounding, or processing of a drug or device that is imported or offered for import into the United States shall, through electronic means in accordance with the criteria of the Secretary-''.}
\uscnotepar{Subsec. (i)(1)(A). Pub. L. 112-144, §~702(b)(1)(B), amended subpar. (A) generally. Prior to amendment, subpar. (A) read as follows: ``upon first engaging in any such activity, immediately register with the Secretary the name and place of business of the establishment, the name of the United States agent for the establishment, the name of each importer of such drug or device in the United States that is known to the establishment, and the name of each person who imports or offers for import such drug or device to the United States for purposes of importation; and''.}
\uscnotepar{Subsec. (i)(1)(B). Pub. L. 112-144, §~702(b)(1)(C), amended subpar. (B) generally. Prior to amendment, subpar. (B) read as follows: ``each establishment subject to the requirements of subparagraph (A) shall thereafter-}
\uscnotepar{``(i) with respect to drugs, register with the Secretary on or before December 31 of each year; and}
\uscnotepar{``(ii) with respect to devices, register with the Secretary during the period beginning on October 1 and ending on December 31 of each year.''}
\uscnotepar{Subsec. (i)(4). Pub. L. 112-144, §~702(b)(2), added par. (4).}
\uscnotepar{Subsec. (j)(1)(E). Pub. L. 112-144, §~703(1), added subpar. (E).}
\uscnotepar{Subsec. (j)(4). Pub. L. 112-144, §~703(2), added par. (4).}
\uscnotepar{Subsec. (n). Pub. L. 112-144, §~604, designated existing provisions as par. (1) and added par. (2).}
\uscnotepar{Subsec. (p). Pub. L. 112-144, §~704, inserted subsec. heading, designated existing provisions as par. (1) and inserted par. heading, and added pars. (2) and (3).}
\uscnotepar{2007-Subsec. (b). Pub. L. 110-85, §~222(a), designated existing provisions as par. (1), struck out ``or a device or devices'' after ``drug or drugs'', and added par. (2).}
\uscnotepar{Subsec. (i)(1). Pub. L. 110-85, §~222(b), inserted text of par. (1) and struck out former text of par. (1) which related to registration requirement for foreign establishments engaged in the manufacture, preparation, propagation, compounding, or processing of a drug or device to be imported or offered for import into the United States.}
\uscnotepar{Subsec. (j)(2). Pub. L. 110-85, §~223, in introductory provisions, substituted ``Each person who registers with the Secretary under this section shall report to the Secretary, with regard to drugs once during the month of June of each year and once during the month of December of each year, and with regard to devices once each year during the period beginning on October 1 and ending on December 31, the following information:'' for ``Each person who registers with the Secretary under this section shall report to the Secretary once during the month of June of each year and once during the month of December of each year the following information:''.}
\uscnotepar{Subsec. (k). Pub. L. 110-85, §~801(b)(3)(C), inserted concluding provisions.}
\uscnotepar{Subsec. (p). Pub. L. 110-85, §~224, amended subsec. (p) generally. Prior to amendment, subsec. (p) read as follows: ``Registrations under subsections (b), (c), (d), and (i) of this section (including the submission of updated information) shall be submitted to the Secretary by electronic means, upon a finding by the Secretary that the electronic receipt of such registrations is feasible, unless the Secretary grants a request for waiver of such requirement because use of electronic means is not reasonable for the person requesting such waiver.''}
\uscnotepar{2004-Subsec. (\textit{o})(1)(B), (2)(B). Pub. L. 108-214, §~2(c)(2)(A), (B)(i), substituted ``or adulterated'' for ``, adulterated''.}
\uscnotepar{Subsec. (\textit{o})(2)(E). Pub. L. 108-214, §~2(c)(2)(B)(ii), substituted ``semi-critical'' for ``semicritical''.}
\uscnotepar{2002-Subsec. (h). Pub. L. 107-250, §~201(e), inserted ``, or by persons accredited to conduct inspections under section 374(g) of this title,'' after ``duly designated by the Secretary''.}
\uscnotepar{Subsec. (i)(1). Pub. L. 107-188, §~321(a)(1), substituted ``On or before December 31 of each year, any establishment'' for ``Any establishment'' and ``shall, through electronic means in accordance with the criteria of the Secretary, register with the Secretary the name and place of business of the establishment, the name of the United States agent for the establishment, the name of each importer of such drug or device in the United States that is known to the establishment, and the name of each person who imports or offers for import such drug or device to the United States for purposes of importation'' for ``shall register with the Secretary the name and place of business of the establishment and the name of the United States agent for the establishment''.}
\uscnotepar{Subsec. (j)(1). Pub. L. 107-188, §~321(a)(2), substituted ``subsection (b), (c), (d), or (i)'' for ``subsection (b), (c), or (d)'' in first sentence.}
\uscnotepar{Subsec. (m)(1). Pub. L. 107-250, §~211, inserted at end ``The Secretary shall publish such list on the Internet site of the Food and Drug Administration. The list so published shall be updated not later than 30 days after each revision of the list by the Secretary.''}
\uscnotepar{Subsec. (\textit{o}). Pub. L. 107-250, §~302(b), added subsec. (\textit{o}).}
\uscnotepar{Subsec. (p). Pub. L. 107-250, §~207, added subsec. (p).}
\uscnotepar{1997-Subsec. (g). Pub. L. 105-115, §~213(b)(3), inserted at end ``In this subsection, the term `wholesale distributor' means any person (other than the manufacturer or the initial importer) who distributes a device from the original place of manufacture to the person who makes the final delivery or sale of the device to the ultimate consumer or user.''}
\uscnotepar{Subsec. (g)(4), (5). Pub. L. 105-115, §~213(b)(1), (2), added par. (4) and redesignated former par. (4) as (5).}
\uscnotepar{Subsec. (i). Pub. L. 105-115, §~417, amended subsec. (i) generally. Prior to amendment, subsec. (i) read as follows: ``Any establishment within any foreign country engaged in the manufacture, preparation, propagation, compounding, or processing of a drug or drugs, or a device or devices, shall be permitted to register under this section pursuant to regulations promulgated by the Secretary. Such regulations shall require such establishment to provide the information required by subsection (j) of this section and shall require such establishment to provide the information required by subsection (j) of this section in the case of a device or devices and shall include provisions for registration of any such establishment upon condition that adequate and effective means are available, by arrangement with the government of such foreign country or otherwise, to enable the Secretary to determine from time to time whether drugs or devices manufactured, prepared, propagated, compounded, or processed in such establishment, if imported or offered for import into the United States, shall be refused admission on any of the grounds set forth in section 381(a) of this title.''}
\uscnotepar{Subsec. (j)(1)(A), (D). Pub. L. 105-115, §~125(a)(2)(C), struck out ``, 356, 357,'' before ``or 360b of this title''.}
\uscnotepar{Subsec. (k). Pub. L. 105-115, §~206(a)(1), inserted ``or person who is accredited under section 360m(a) of this title'' after ``report to the Secretary''.}
\uscnotepar{Subsecs. (\textit{l}), (m). Pub. L. 105-115, §~206(a)(2), added subsecs. (\textit{l}) and (m).}
\uscnotepar{Subsec. (n). Pub. L. 105-115, §~209(a), added subsec. (n).}
\uscnotepar{1976-Subsec. (a)(1). Pub. L. 94-295, §~4(a)(2), substituted ``drug package or device package'' for ``drug package'', ``distribution of the drug or device'' for ``distribution of the drug'', and ``ultimate consumer or user'' for ``ultimate consumer''.}
\uscnotepar{Subsecs. (b) to (d). Pub. L. 94-295, §~4(a)(3), inserted ``or a device or devices'' after ``drug or drugs''.}
\uscnotepar{Subsec. (e). Pub. L. 94-295, §~4(a)(4), authorized the Secretary to prescribe by regulation a uniform system for the identification of devices intended for human use and authorized him, in addition, to require that persons who are required to list devices pursuant to subsec. (j) also list such devices in accordance with the system.}
\uscnotepar{Subsec. (g)(1) to (3). Pub. L. 94-295, §~4(a)(5), substituted ``drugs or devices'' for ``drugs''.}
\uscnotepar{Subsec. (h). Pub. L. 94-295, §~4(a)(6), inserted reference to establishments engaged in the manufacture, propagation, compounding, or processing of a drug or drugs or of a device or devices classified in class II or III.}
\uscnotepar{Subsec. (i). Pub. L. 94-295, §~4(a)(7), inserted reference to devices and inserted requirement that regulations require establishments to provide the information required by subsection (j) of this section in the case of a device or devices.}
\uscnotepar{Subsec. (j)(1). Pub. L. 94-295, §~4(a)(8)(A), in introductory provisions substituted ``a list of all drugs and a list of all devices and a brief statement of the basis for believing that each device included in the list is a device rather than a drug (with each drug and device in each list listed by its established name'' for ``a list of all drugs (by established name'' and ``drugs or devices filed'' for ``drugs filed''.}
\uscnotepar{Subsec. (j)(1)(A). Pub. L. 94-295, §~4(a)(8)(B), substituted ``the applicable list'' for ``such list'', inserted ``or a device intended for human use contained in the applicable list with respect to which a performance standard has been established under section 360d of this title or which is subject to section 360e of this title,'' after ``360b of this title,'', and substituted ``such drug or device'' for ``such drug'' wherever appearing.}
\uscnotepar{Subsec. (j)(1)(B). Pub. L. 94-295, §~4(a)(8)(C), in introductory provisions substituted ``drug or device contained in an applicable list'' for ``drug contained in such list''.}
\uscnotepar{Subsec. (j)(1)(B)(i). Pub. L. 94-295, §~4(a)(8)(D), substituted ``which drug is subject to section 353(b)(1) of this title, or which device is a restricted device, a copy of all labeling for such drug or device, a representative sampling of advertisements for such drug or device, and, upon request made by the Secretary for good cause, a copy of all advertisements for a particular drug product or device, or'' for ``which is subject to section 353(b)(1) of this title, a copy of all labeling for such drug, a representative sampling of advertisements for such drug, and, upon request made by the Secretary for good cause, a copy of all advertisements for a particular drug product, or''.}
\uscnotepar{Subsec. (j)(1)(B)(ii). Pub. L. 94-295, §~4(a)(8)(E), substituted ``which drug is not subject to section 353(b)(1) of this title or which device is not a restricted device, the label and package insert for such drug or device and a representative sampling of any other labeling for such drug or device'' for ``which is not subject to section 353(b)(1) of this title, the label and package insert for such drug and a representative sampling of any other labeling for such drug''.}
\uscnotepar{Subsec. (j)(1)(C). Pub. L. 94-295, §~4(a)(8)(F), substituted ``an applicable list'' for ``such list''.}
\uscnotepar{Subsec. (j)(1)(D). Pub. L. 94-295, §~4(a)(8)(G), substituted ``a list'' for ``the list'', inserted ``or the particular device contained in such list is not subject to a performance standard established under section 360d of this title or to section 360e of this title or is not a restricted device'' after ``or 360b of this title,'', and substituted ``particular drug product or device'' for ``particular drug product'' wherever appearing.}
\uscnotepar{Subsec. (j)(2). Pub. L. 94-295, §~4(a)(8)(H), substituted ``drug or device'' for ``drug'' in subpars. (A), (B), and (C), and substituted ``(each by established name'' for ``(by established name'' in subpar. (C).}
\uscnotepar{Subsec. (k). Pub. L. 94-295, §~4(a)(9), added subsec. (k).}
\uscnotepar{1972-Subsec. (e). Pub. L. 92-387, §~4(a), inserted provision that the Secretary may assign a listing number to each drug or class of drugs listed under subsec. (j).}
\uscnotepar{Subsec. (f). Pub. L. 92-387, §~4(b), inserted exception that the list submitted under subsec. (j)(3) and information submitted under subsec. (j)(1), (2) shall be exempt from inspection unless the Secretary determines otherwise.}
\uscnotepar{Subsec. (i). Pub. L. 92-387, §~4(c), inserted provision that the regulations shall require such establishment to provide the information required by subsec. (j).}
\uscnotepar{Subsec. (j). Pub. L. 92-387, §~3, added subsec. (j).}
\uscnotepar{1970-Subsec. (a). Pub. L. 91-513 struck out provisions defining the wholesaling, jobbing, or distributing of depressant or stimulant drugs.}
\uscnotepar{Subsec. (b). Pub. L. 91-513 struck out provisions covering establishments engaged in the wholesaling, jobbing, or distributing of depressant or stimulant drugs and the inclusion of the fact of such activity in the annual registration.}
\uscnotepar{Subsec. (c). Pub. L. 91-513 struck out provisions covering new registrations of persons first engaging in the wholesaling, jobbing, or distributing of depressant or stimulant drugs and the inclusion of the fact of such activity in the registration.}
\uscnotepar{Subsec. (d). Pub. L. 91-513 struck out number designation ``(1)'' preceding first sentence, struck out portion of such redesignated provisions covering the wholesaling, jobbing, or distributing of depressant or stimulant drugs, and struck out par. (2) covering the filing of supplemental registration whenever a person not previously engaged or involved with depressant or stimulant drugs goes into the manufacturing, preparation, or processing thereof.}
\uscnotepar{1965-Pub. L. 89-74, §~4(e), included certain wholesalers in section catchline.}
\uscnotepar{Subsec. (a)(2), (3). Pub. L. 89-74, §~4(a), added par. (2) and redesignated former par. (2) as (3).}
\uscnotepar{Subsecs. (b), (c). Pub. L. 89-74, §~4(b), (c), inserted ``or in the wholesaling, jobbing, or distributing of any depressant or stimulant drug'' after ``drug or drugs'' and inserted requirement that establishment indicate activity in depressant or stimulant drugs at time of registration.}
\uscnotepar{Subsec. (d). Pub. L. 89-74 §~4(d), designated existing provisions as par. (1), inserted ``or the wholesaling, jobbing, or distributing of any depressant or stimulant drug'' and the requirement that the additional establishment indicate activity in depressant or stimulant drugs at time of registration, and added par. (2).}
\usccrosshead{Statutory Notes and Related Subsidiaries}
\uscnotehead{Effective Date of 2020 Amendment}
\uscnotepar{Amendment by Pub. L. 116-136 effective 180 days after Mar. 27, 2020, see section 3112(g) of Pub. L. 116-136, set out as a note under section 356c of this title.}
\uscnotehead{Effective Date of 2002 Amendment}
\uscnotepar{Amendment by Pub. L. 107-188 effective upon the expiration of the 180-day period beginning June 12, 2002, see section 321(c) of Pub. L. 107-188, set out as a note under section 331 of this title.}
\uscnotehead{Effective Date of 1997 Amendment}
\uscnotepar{Amendment by sections 206(a), 209(a), 213(b), and 417 of Pub. L. 105-115 effective 90 days after Nov. 21, 1997, except as otherwise provided, see section 501 of Pub. L. 105-115, set out as a note under section 321 of this title.}
\uscnotehead{Effective Date of 1972 Amendment}
\uscnotepar{Pub. L. 92-387, §~5, Aug. 16, 1972, 86 Stat. 562, provided that: ``The amendments made by this Act [amending this section and sections 331 and 335 of this title and enacting provisions set out below] shall take effect on the first day of the sixth month beginning after the date of enactment of this Act [Aug. 16, 1972].''}
\uscnotehead{Effective Date of 1970 Amendment}
\uscnotepar{Amendment by Pub. L. 91-513 effective on first day of seventh calendar month that begins after Oct. 26, 1970, see section 704 of Pub. L. 91-513, set out as an Effective Date note under section 801 of this title.}
\uscnotehead{Effective Date of 1965 Amendment}
\uscnotepar{Amendment by Pub. L. 89-74 effective Feb. 1, 1966, subject to registration with Secretary of names, places of business, establishments, and other prescribed information prior to Feb. 1, 1966, see section 11 of Pub. L. 89-74, set out as a note under section 321 of this title.}
\uscnotehead{Updating Regulations}
\uscnotepar{Pub. L. 117-328, div. FF, title II, §~2511(b), Dec. 29, 2022, 136 Stat. 5804, provided that: ``Not later than 2 years after the date of enactment of this Act [Dec. 29, 2022], the Secretary of Health and Human Services shall update regulations, as appropriate, to implement the amendment made by subsection (a) [amending this section].''}
\uscnotehead{Savings Provision}
\uscnotepar{Amendment by Pub. L. 91-513 not to affect or abate any prosecutions for any violation of law or any civil seizures or forfeitures and injunctive proceedings commenced prior to the effective date of such amendment, and all administrative proceedings pending before the Bureau of Narcotics and Dangerous Drugs [now the Drug Enforcement Administration] on Oct. 27, 1970, to be continued and brought to final determination in accord with laws and regulations in effect prior to Oct. 27, 1970, see section 702 of Pub. L. 91-513, set out as a note under section 321 of this title.}
\uscnotehead{Device Modifications}
\uscnotepar{Pub. L. 114-255, div. A, title III, §~3059(b), Dec. 13, 2016, 130 Stat. 1130, provided that: ``The Secretary of Health and Human Services, acting through the Commissioner of Food and Drugs, shall issue final guidance regarding when a premarket notification under section 510(k) of the Federal Food, Drug, and Cosmetic Act (21 U.S.C. 360(k)) is required to be submitted for a modification or change to a legally marketed device. Such final guidance shall be issued not later than 1 year after the date on which the comment period closes for the draft guidance on such subject.''}
\uscnotehead{Declaration of Policy of Drug Listing Act of 1972}
\uscnotepar{Pub. L. 92-387, §~2, Aug. 16, 1972, 86 Stat. 559, provided that: ``The Federal Government which is responsible for regulating drugs has no ready means of determining what drugs are actually being manufactured or packed by establishments registered under the Federal Food, Drug, and Cosmetic Act [21 U.S.C. 301 et seq.] except by periodic inspection of such registered establishments. Knowledge of which particular drugs are being manufactured or packed by each registered establishment would substantially assist in the enforcement of Federal laws requiring that such drugs be pure, safe, effective, and properly labeled. Information on the discontinuance of a particular drug could serve to alleviate the burden of reviewing and implementing enforcement actions against drugs which, although commercially discontinued, remain active for regulatory purposes. Information on the type and number of different drugs being manufactured or packed by drug establishments could permit more effective and timely regulation by the agencies of the Federal Government responsible for regulating drugs, including identification of which drugs in interstate commerce are subject to section 505 or 507 [21 U.S.C. 355, 357], or to other provisions of the Federal Food, Drug, and Cosmetic Act.''}
\uscnotehead{Congressional Declaration of Need for Registration and Inspection of Drug Establishments}
\uscnotepar{Pub. L. 87-781, title III, §~301, Oct. 10, 1962, 76 Stat. 793, provided that: ``The Congress hereby finds and declares that in order to make regulation of interstate commerce in drugs effective, it is necessary to provide for registration and inspection of all establishments in which drugs are manufactured, prepared, propagated, compounded, or processed; that the products of all such establishments are likely to enter the channels of interstate commerce and directly affect such commerce; and that the regulation of interstate commerce in drugs without provision for registration and inspection of establishments that may be engaged only in intrastate commerce in such drugs would discriminate against and depress interstate commerce in such drugs, and adversely burden, obstruct, and affect such interstate commerce.''}
\uscnotehead{Registration of Certain Persons Owning or Operating Drug Establishments Prior to Oct. 10, 1962}
\uscnotepar{Pub. L. 87-781, title III, §~303, Oct. 10, 1962, 76 Stat. 795, provided that any person who, on the day immediately preceding Oct. 10, 1962, owned or operated an establishment which manufactured or processed drugs, registered before the first day of the seventh month following October, 1962, would be deemed to be registered in accordance with subsec. (b) of this section for the calendar year 1962 and if registered within this period and effected in 1963, be deemed in compliance for that calendar year.}
\usccrosshead{Executive Documents}
\uscnotehead{Ex. Ord. No. 14293. Regulatory Relief To Promote Domestic Production of Critical Medicines}
\uscnotepar{Ex. Ord. No. 14293, May 5, 2025, 90 F.R. 19615, provided:}
\uscnotepar{By the authority vested in me as President by the Constitution and the laws of the United States of America, it is hereby ordered:}
\uscnotepar{Section 1. \textit{Purpose}. During my first term, my Administration took unprecedented action to improve the well-being of the American people by restoring capacity for domestic production of critical pharmaceutical products. Notably, in Executive Order 13944 of August 6, 2020 (Combating Public Health Emergencies and Strengthening National Security By Ensuring Essential Medicines, Medical Countermeasures, and Critical Inputs Are Made In The United States) [42 U.S.C. 247d-6b note], I directed each executive department and agency involved in the procurement of Essential Medicines, Medical Countermeasures, and Critical Inputs to take a variety of actions to increase their domestic procurement of Essential Medicines, Medical Countermeasures, and Critical Inputs, as defined in section 7 of that order, and to identify vulnerabilities in our Nation's supply chains for these products. Unfortunately, the prior administration did too little to advance these goals. Critical barriers and information gaps persist in establishing a domestic, resilient, and affordable pharmaceutical supply chain for American patients.}
\uscnotepar{One key area of concern is the length of time it takes to build pharmaceutical manufacturing facilities in the United States today. New construction must navigate myriad Federal, State, and local requirements ranging from building standards and zoning restrictions to environmental protocols that together diminish the certainty needed to generate investment for large manufacturing projects. For pharmaceutical manufacturing, these barriers are heightened by unannounced inspections of domestic manufacturers by the Food and Drug Administration (FDA), which are more frequent than such inspections at international facilities. Industry estimates suggest that building new manufacturing capacity for pharmaceuticals and critical inputs may take as long as 5 to 10 years, which is unacceptable from a national security standpoint. Even expanding existing capacity or modifying existing production lines to produce new or different products requires extensive permitting and regulatory approval, making it more difficult to repurpose existing underutilized pharmaceutical manufacturing capacity available domestically.}
\uscnotepar{It is in the best interest of the Nation to eliminate regulatory barriers to the domestic production of the medicines Americans need. My Administration will work to make the United States the most competitive nation in the world for the manufacture of safe and effective pharmaceutical products.}
\uscnotepar{Sec. 2. \textit{Policy}. It is the policy of the United States that the regulation of manufacturing pharmaceutical products and inputs be streamlined to facilitate the restoration of a robust domestic pharmaceutical manufacturing base.}
\uscnotepar{Sec. 3. \textit{Streamlining Review of Domestic Pharmaceutical Manufacturing by the Food and Drug Administration}. Within 180 days of the date of this order [May 5, 2025], the Secretary of Health and Human Services, through the Commissioner of Food and Drugs (FDA Commissioner), shall review existing regulations and guidance that pertain to the development of domestic pharmaceutical manufacturing and shall take steps to eliminate any duplicative or unnecessary requirements in such regulations and guidance; maximize the timeliness and predictability of agency review; and streamline and accelerate the development of domestic pharmaceutical manufacturing. The FDA Commissioner's review shall encompass all regulations and guidance that apply to the inspection and approval of new and expanded manufacturing capacity, emerging technologies that enable the manufacturing of pharmaceutical products, active pharmaceutical ingredients, key starting materials, and associated raw materials in the United States. The FDA Commissioner shall:}
\uscnotepar{(a) evaluate the current risk-based approach to prior approval of licensure inspections, including when such inspections are necessary, and seek to improve upon this approach to ensure all required inspections are prompt, efficient, and limited to what is necessary to ensure compliance with the Federal Food, Drug, and Cosmetic Act [21 U.S.C. 301 et seq.] and other Federal law;}
\uscnotepar{(b) identify and undertake measures necessary to expand, as practicable, existing programs that provide early technical advice before a facility is operational}
\uscnotepar{;(c) identify and undertake measures necessary to improve enforcement of data reporting under section 510(j)(3) of the Federal Food, Drug, and Cosmetic Act (21 U.S.C. 360(j)(3)), including consideration of publicly displaying the list of facilities, including foreign facilities, that are not in compliance;}
\uscnotepar{(d) provide clearer guidance regarding the requirements or recommendations for site changes, including moving production from a foreign to domestic facility, and validation of new or updated components necessary in manufacturing; and}
\uscnotepar{(e) review and, as appropriate, seek to update any other relevant compliance policies, guidance documents, and regulations.}
\uscnotepar{Sec. 4. \textit{Enhancing Inspection of Foreign Manufacturing Facilities}. Within 90 days of the date of this order, the FDA Commissioner shall develop and advance improvements to the risk-based inspection regime that ensures routine reviews of overseas manufacturing facilities involved in the supply of United States medicines, which shall be funded by increased fees on foreign manufacturing facilities to the extent consistent with applicable law. Additionally, the FDA Commissioner shall publicly disclose the annual number of inspections that the FDA conducts on such foreign facilities, with specific detail by country and by manufacturer.}
\uscnotepar{Sec. 5. \textit{Streamlining Review of Domestic Pharmaceutical Manufacturing by the Environmental Protection Agency}. Within 180 days of the date of this order, the Administrator of the Environmental Protection Agency (EPA) shall take action to update regulations and guidance that apply to the inspection and approval of new and expanded manufacturing capacity of pharmaceutical products, active pharmaceutical ingredients, key starting materials, and associated raw materials in the United States to eliminate any duplicative or unnecessary requirements and maximize the timeliness and predictability of agency review.}
\uscnotepar{Sec. 6. \textit{Centralized Coordination of Environmental Permits to Expand Domestic Pharmaceutical Manufacturing Capacity}. For purposes of 42 U.S.C. 4336a [section 107 of the National Environmental Policy Act of 1969], the EPA shall be the lead agency for the permitting of pharmaceutical manufacturing facilities that require preparation of an Environmental Impact Statement pursuant to the National Environmental Policy Act of 1969, 42 U.S.C. 4321 et seq., unless that role is assumed by another agency. The lead agency shall designate a single point of contact within the agency to coordinate with permit applicants. The Office of Management and Budget shall coordinate with the lead agency and with other relevant agencies and the Federal Permitting Improvement Steering Committee, as needed, to expedite the review and approval of relevant permits.}
\uscnotepar{Sec. 7. \textit{Streamlining Review of Domestic Pharmaceutical Manufacturing by the United States Army Corps of Engineers}. Within 180 days of the date of this order, the Secretary of the Army, acting through the Assistant Secretary of the Army for Civil Works, shall review the nationwide permits issued under section 404 of the Clean Water Act of 1972 [probably means section 404 of the Federal Water Pollution Control Act] (33 U.S.C. 1344) and section 10 of the Rivers and Harbors Appropriation Act of 1899 (33 U.S.C. 403) to determine whether an activity-specific nationwide permit is needed to facilitate the efficient permitting of pharmaceutical manufacturing facilities.}
\uscnotepar{Sec. 8. \textit{General Provisions}. (a) Nothing in this order shall be construed to impair or otherwise affect:}
\uscnotepar{(i) the authority granted by law to an executive department or agency, or the head thereof; or}
\uscnotepar{(ii) the functions of the Director of the Office of Management and Budget relating to budgetary, administrative, or legislative proposals.}
\uscnotepar{(b) This order shall be implemented consistent with applicable law and subject to the availability of appropriations.}
\uscnotepar{(c) This order is not intended to, and does not, create any right or benefit, substantive or procedural, enforceable at law or in equity by any party against the United States, its departments, agencies, or entities, its officers, employees, or agents, or any other person.}
\uscnotepar{(d) The Department of Health and Human Services shall provide funding for publication of this order in the Federal Register.}
\uscnotepar{Donald J. Trump.}

%% - DRAFTING INSTRUCTIONS (added for VVUQ-04 draft-bill; the reproduced statutory text above is unchanged) -
\begin{draftbox}{§ 360 (21 U.S.C. § 360) - registration of producers; the 510(k) premarket notification in subsection (k)}
\dstep{Role in the amendment: subsection (k) of § 360 is the 510(k) premarket-notification pathway. The conforming amendment harmonizes this pathway with the verification protocol in new section 515D and the predetermined change control plan authority in section 360e-4, so a verification-governed change to a Physical AI device does not require a new premarket notification.}
\dstep{Critical naming: use § 360(k) (subsection (k) of section 360) for the 510(k) premarket notification, and reserve § 360k (the separate section) for state preemption handled in sections/s360k.tex. Do not conflate the two; this is the checklist correction in papers/VVUQ-04/instruct-bill/10-legal-crosswalk-and-bill-style.md (Section E, item 3).}
\dstep{Amendatory action: Section 510 of the Federal Food, Drug, and Cosmetic Act (21 U.S.C. § 360) is amended in subsection (k) by adding at the end the following: "A change to a device that generates or executes robot-patient interaction code, cleared under this subsection, that is consistent with a verification protocol under section 515D and an established predetermined change control plan under section 360e-4, shall not require a new premarket notification under this subsection." Confirm consistency with § 360e-4(b), which already excuses notification for changes consistent with an established plan.}
\dstep{Source synthesis: process papers/VVUQ-04/instruct-bill/01-federal-statutory-framework.md (the 510(k) pathway, expressly distinct from § 360k) and 02-fda-ai-device-regulation.md (the PCCP final guidance and the 510(k) change pathway) together with papers/VVUQ-03/final-paper/sections/statutory\_text.tex. Resolve to references.bib keys usc-fdca, usc-360e4, and fda-pccp-2024.}
\dstep{Comparative print rendering: reproduce subsection (k) and the surrounding subsections unchanged, and show the added sentence wrapped in \cprinsert{...} at the end of (k); leave subsections (a) through (j) and (l) through (q) unchanged.}
\dstep{Formatting: single hyphens only; § for every codified reference; even RaggedRight spacing; reword so the inserted sentence does not end in a one- or two-word line.}
\end{draftbox}
